% TRANSLATION-PRODUCER DRAFT — UNCHECKED
% Work: Emmy Noether, Paper 41 — Korean translation unit U05
% Source custody snapshot: ${PUBLIC_INTERLANGUAGE_ROOT}\03_projects\language_management\cjk\03_working_translations\noether_paper41_zh_translation_001_20260722\source\Noether_Paper41_CurrentGermanAuthority_interval.tex
% Source snapshot lines: 58--66 (historical whole-source lines 19766--19774)
% Source unit: 2,671 LF-normalized UTF-8 bytes; SHA-256 BCEA8450C5DC412DA1506DD0A8A935AD5517D2B03637C57698B6C3D504FB7788
% Source snapshot SHA-256: C265058425E5E2D1A2289CC03A9DDEDDDF4803A3215DC3F173B93E7AB69D60ED
% Historical whole-source SHA-256: 443EF950D7D45DC6D9E44A9B87501620C10DA873E50E5F2B253ECCAE6A946D27
% Producer role: translation only. No source/scan check, Korean review, compilation, rendering, assembly, packaging, certification, approval, or self-check was performed.
% Handoff state: UNCHECKED; requires an independent Korean checker.

\noindent\textbf{2.} \emph{최소형 주종정리}의 정식화에서는 관계 (2)부터 (5)가 순전히 곱셈적이라는 사실을 사용한다. 따라서 이 관계들은 동시에 \(n\)개의 잉여류 \(u_SK^*\)의 원소들로 이루어진 \(\mathfrak G\)의 확장군 \(\mathfrak G^*\)을 정의한다:

\srcdisplay{(6)}{\mathfrak G^*=\{u_{S_1}K^*,\ldots,u_{S_n}K^*\}.}

\(\mathfrak G^*\)은 \(K^*\)을 아벨 정규부분군으로 가지며, \(\mathfrak G^*/K^*\)은 \(\mathfrak G\)와 동형이다\srcfn{8)}{따라서 관계 (2)부터 (5)는 단순히 군의 확장에 관한 정의관계이다. 이러한 확장은, \(K^*\)을 오직 곱셈군으로 보았을 때 그 전체 자기동형군 안에 \(\mathfrak G\)의 준동형상인 부분군이 존재하기만 하면 성립한다는 것이 알려져 있다. 교차곱의 특수성은 이 준동형이 동형이 되고, 더 나아가 바로 그 부분군으로 곱셈적이면서 동시에 덧셈적인 자기동형사상 전체, 즉 \(K\)의 갈루아군을 택한다는 데 있다. 이로써 군의 확장과 더불어 환의 확장도 얻어진다.}. \(\mathfrak G^*\)은 \(K\)를 자기 자신으로 변환하는 모든 (정칙) 원소 \(g^*\), 즉 \(g^{*-1}Kg^*=K\)인 원소들로 이루어진다. 그러므로 \(K\)의 환 자기동형사상은 정확히 \(\mathfrak G^*\)의 원소들에 의한 변환으로 유도된다. 실제로 그러한 \(g^*\)는 모두 \(K\)의 환 자기동형사상을 유도한다. 역으로 그러한 자기동형사상은 모두 어떤 원소 \(v=u_Sw\)에 의한 변환으로 유도되는데, 여기서 \(w\)는 \(K\) 위에 항등사상을 유도한다. 따라서 \(w\)는 \(K\)의 모든 원소와 가환하므로 \(K\)에 속한다. \(K\)가 최대 가환 부분환이므로 \(w\)는 \(K^*\)에 속한다.

\emph{최소형 주종정리. 제1정식: \(K^*\)의 항등자기동형사상의 연장인 \(\mathfrak G^*\)의 모든 군 자기동형사상은 내부자기동형사상이며, \(K^*\)의 한 원소에 의해 유도된다. 제2정식: 변환량 \(c_S^Tc_T/c_{ST}\)가 모두 1이면, \(c_S\)들은 기호적 \((1-S)\)제곱이다. 즉, 모든 \(S\in\mathfrak G\)에 대하여 \(c_S=b^{1-S}=b/b^S\)가 되게 하는 \(b\in K^*\)가 존재한다. 제3정식: 군 \(\mathfrak G\)은 \(K^*\) 안에서 인자계 1에 속하는 1차 교차표현류를 단 하나만 가진다.}

이때 표현 \(u_S\mapsto C_S\)가 인자계 \(a_{S,T}\)에 대한 교차표현이라는 것은 \(C_S^TC_T=C_{ST}a_{S,T}\)가 성립한다는 뜻이다. 두 표현은 \(C_S=B^{-S}D_SB\)일 때 같은 류에 속한다.
